\documentclass[11pt]{amsart}
\usepackage[margin=1in]{geometry}
\usepackage{lmodern,microtype,amsmath,amssymb,amsthm}
\usepackage[hidelinks]{hyperref}
\newtheorem{theorem}{Theorem}[section]
\newtheorem{proposition}[theorem]{Proposition}
\newtheorem{lemma}[theorem]{Lemma}
\newtheorem{corollary}[theorem]{Corollary}
\theoremstyle{definition}
\newtheorem{definition}[theorem]{Definition}
\newtheorem{remark}[theorem]{Remark}
\newcommand{\R}{\mathbb R}
\DeclareMathOperator{\supp}{supp}
\DeclareMathOperator{\Vol}{Vol}
\emergencystretch=2em
\title{Isometric tight tori with different enclosed volumes}
\author{Yuetong Zhang}
\address{Peking University}
\email{2300010810@stu.pku.edu.cn}
\email{gzhangyuetong@gmail.com}
\date{October 31, 2026}
\hypersetup{pdftitle={Isometric tight tori with different enclosed volumes},pdfauthor={Yuetong Zhang}}
\begin{document}
\begin{abstract}
We construct two smooth embedded tight tori in Euclidean three-space
with identical induced metrics and different enclosed volumes.
Starting with a closed asymptotic curve in the annulus of Banks,
Ghomi, Qiu, and Raffaelli, we complete the annulus to a tight torus.
More generally, any tight torus containing a closed asymptotic curve
in negative curvature can be modified near the curve to admit a
compactly supported infinitesimal bending with nonzero normal flux.
The key steps are a periodic ruled-annulus construction and a
Codazzi-based bending formula. Nonzero flux follows from a
differentiated area identity. Opposite small multiples of the
bending produce the required pair.
\end{abstract}
\maketitle

\section{Introduction}

The main result follows from the geometric constructions below.
Conventions are collected in Section~\ref{sec:preliminaries};
all maps and fields are smooth.

\begin{theorem}[Main result]\label{thm:main}
There exist a smooth tight embedding $X:T^2\to\R^3$ and an
infinitesimal bending $Y:T^2\to\R^3$ such that
\[
 \supp Y\Subset\{K_X<0\},
 \qquad \int_{T^2}\langle Y,\nu_X\rangle\,dA_X\ne0.
\]
There is $a_0>0$ such that for every $0<a<a_0$ the maps
$X_\pm=X\pm aY$ are smooth tight embeddings satisfying
\begin{equation}\label{eq:main-pair}
 g_{X_+}=g_{X_-},\qquad \Vol(X_+)\ne\Vol(X_-).
\end{equation}
Consequently their images are noncongruent. Both maps agree with $X$
outside $\supp Y$, and in particular on a nonempty open set.
\end{theorem}

The construction starts with the annulus of Banks, Ghomi, Qiu, and
Raffaelli and ends with opposite multiples of one supported
infinitesimal bending.

\begin{theorem}[Existence of a starting torus]\label{thm:seed}
There exist a smooth tight embedding $X_0:T^2\to\R^3$ and a smooth
embedding $c:S^1\to T^2$ such that $X_0\circ c$ is asymptotic and
$K_{X_0}<0$ on $c(S^1)$.
\end{theorem}

\emph{Method.} Flatten the ends of the BGQR annulus to infinite order
and attach it to a convex sphere along planar caps. The retained
asymptotic curve lies in negative curvature, and total positive
curvature $4\pi$ gives tightness (Section~\ref{sec:seed}).

\begin{theorem}[One asymptotic curve gives a prepared core]
\label{thm:prepared-core}
Let $X_0:T^2\to\R^3$ be a tight embedding and let $c:S^1\to T^2$
be an embedded asymptotic curve in $\{K_{X_0}<0\}$. For every open
$U$ with $c(S^1)\subset U\subset\{K_{X_0}<0\}$ and every
$\varepsilon>0$, there is a tight embedding $X_1:T^2\to\R^3$ such that
\[
 \supp(X_1-X_0)\Subset U,\qquad
 \|X_1-X_0\|_{C^2}<\varepsilon,\qquad K_{X_1}<0\text{ on }U.
\]
The curve $X_1\circ c$ remains asymptotic. There are $\ell>0$ and a
reparametrization $c_1:\R/\ell\mathbb Z\to T^2$ of $c$ such that
$\gamma=X_1\circ c_1$ is unit speed. Put $n=\nu_{X_1}\circ c_1$
and $\rho=(-K_{X_1}\circ c_1)^{-1/4}$. At some $s_0$,
\[
 \gamma_{ss}(s_0)=0,\qquad (\rho n)_{ss}(s_0)\ne0.
\]
\end{theorem}

\emph{Method.} Injectivity of the negative Gauss map, spherical
Gauss--Bonnet, and Fenchel's theorem force an inflection. A local
change of speed, corrected to preserve closedness, makes the
curvature-weighted normal accelerate there
(Sections~\ref{sec:frame}--\ref{sec:preparation}).

\begin{theorem}[A prepared core gives periodic ruled annuli]
\label{thm:ruled-interface}
Let $\ell>0$ and let $\gamma:\R/\ell\mathbb Z\to\R^3$ be a
unit-speed embedded asymptotic curve in negative curvature on an
oriented immersed surface. Let $n$ be its surface unit normal and put
$\rho=(-K)^{-1/4}$ along the curve. Suppose that at some $s_0$,
\[
 \gamma_{ss}(s_0)=0,\qquad (\rho n)_{ss}(s_0)\ne0.
\]
There are a smooth periodic field $V$ and $a>0$ such that
$R(s,u)=\gamma(s)+uV(s)$ is an embedded negatively curved strip on
$|u|<a$, with tangent plane $n(s)^\perp$ along the core.
For every open neighborhood $N$ of the core in $\R^3$, there are a
nonempty bounded open interval $I$ and an embedded negatively curved annulus
$B:(\R/\ell\mathbb Z)\times I\to\R^3$ whose image lies in
$N\cap R((\R/\ell\mathbb Z)\times(-a,a))$, whose constant-$c$
circles are asymptotic, and whose constant-$s$ curves are straight.
\end{theorem}

\emph{Method.} Choose the transverse ruling so that the asymptotic
return map is the identity. Concentrating the period correction near
the inflection keeps the ruling transverse
(Section~\ref{sec:ruled}).

\begin{theorem}[Replacement along an asymptotic core]
\label{prop:insert}
Let $\ell>0$, and let $X_0:M\to\R^3$ be a smooth embedding of a closed
oriented surface. Let $c:\R/\ell\mathbb Z\to M$ be an embedded asymptotic
curve with image $\gamma=X_0\circ c$, and let $U$ be an open
neighborhood of $c(\R/\ell\mathbb Z)$ contained in
$\{K_{X_0}<0\}$.
Suppose $b>0$ and
$Z:(\R/\ell\mathbb Z)\times(-b,b)\to\R^3$ is a smooth embedding
such that $Z(s,0)=\gamma(s)$ and the tangent plane spanned by
$Z_s(s,0),Z_u(s,0)$ equals $dX_0(T_{c(s)}M)$.
For every $\varepsilon>0$ there is a smooth embedding
$\widetilde X:M\to\R^3$ such that
\[
 \supp(\widetilde X-X_0)\Subset U,\qquad
 \|\widetilde X-X_0\|_{C^1}<\varepsilon,\qquad
 K_{\widetilde X}<0\text{ on }U,
\]
and there exist $0<a<b$ and a smooth embedding
$q:(\R/\ell\mathbb Z)\times(-a,a)\to U$ with
\[
 q(s,0)=c(s),\qquad \widetilde X(q(s,u))=Z(s,u).
\]
If $X_0$ is tight, $\widetilde X$ is tight.
\end{theorem}

\emph{Method.} Express the new strip as a quadratic normal
displacement and cut it off in a narrow collar. An estimate of the
second fundamental form preserves negative curvature through the
transition (Section~\ref{sec:insertion}).

\begin{theorem}[Geometric band theorem]\label{thm:band-bending}
Let $\ell>0$, let $I\subset\R$ be a nonempty open interval, and let
\[
 B:(\R/\ell\mathbb Z)\times I\longrightarrow\R^3
\]
be a smooth embedding with $K_B<0$. Suppose each constant-$c$
circle is asymptotic, and each constant-$s$ curve lies on a straight
line. Then, for every nonempty open subinterval $I_0\subset I$,
there is an infinitesimal bending $Y$ of $B$ such that
\[
 \supp Y\Subset(\R/\ell\mathbb Z)\times I_0,\qquad
 \int_{(\R/\ell\mathbb Z)\times I}
                \det(Y,B_s,B_c)\,ds\,dc\ne0.
\]
\end{theorem}

\emph{Method.} The Codazzi identities give a family of supported
bendings; differentiated area identities show that its flux
functional is nonzero (Section~\ref{sec:bending}). On an inserted
band, extension by zero gives a global bending
(Lemma~\ref{lem:globalize}).

\begin{theorem}[A supported bending implies an isometric pair]
\label{prop:pair}
Suppose $X:M\to\R^3$ is a smooth tight embedding of a closed surface,
and $Y$ is an infinitesimal bending with
\[
 \supp Y\Subset\{K_X<0\},\qquad
 \int_M\langle Y,\nu_X\rangle\,dA_X\ne0.
\]
Then, for all sufficiently small $a>0$, $X_\pm=X\pm aY$ are tight
embeddings with unequal enclosed volumes and
\begin{equation}\label{eq:exact-metric}
 g_{X_\pm}=g_X+a^2\langle dY,dY\rangle,
 \quad
 \langle dY,dY\rangle(v,w)=\langle dY(v),dY(w)\rangle.
\end{equation}
\end{theorem}

\emph{Method.} The bending equation cancels the linear metric term,
while nonzero flux gives a nonzero first volume variation.
Small amplitudes preserve embeddedness and tightness
(Section~\ref{sec:reductions}). The volume-stationarity problem is
discussed by V.~A.~Aleksandrov \cite{Aleksandrov}.

\begin{proof}[Proof of Theorem~\ref{thm:main} from the roadmap]
Theorem~\ref{thm:seed} gives the starting torus and curve.
The preparation, ruling, insertion, and band theorems yield the
global supported bending, as assembled in
Corollary~\ref{thm:one-curve}. Apply Theorem~\ref{prop:pair}.
The positive-curvature locus is nonempty by
Theorem~\ref{thm:tightness} and remains fixed, giving the common
open set.
\end{proof}

\section{Preliminaries}\label{sec:preliminaries}

All objects are $C^\infty$. A closed surface is compact, connected,
and without boundary, and $T^2=(\R/2\pi\mathbb Z)^2$.
For an immersion $X$, write $g_X=\langle dX,dX\rangle$, $K_X$,
and $dA_X$ for its induced metric, Gaussian curvature, and area
measure. Our sign convention is
$\mathrm{II}_{ij}=\langle X_{ij},n\rangle$ for a chosen unit normal.
The Gauss equation is
\begin{equation}\label{eq:gauss-curvature}
 K_X=\frac{\mathrm{II}_{11}\mathrm{II}_{22}
                 -\mathrm{II}_{12}^2}{g_{11}g_{22}-g_{12}^2}.
\end{equation}
See \cite[\S3-3 and \S4-3]{doCarmo}. We use the Levi--Civita
connection $\nabla$ and its Christoffel coefficients $\Gamma^k_{ij}$;
repeated coordinate indices are summed.

\begin{definition}\label{def:tight}
An embedding of a closed surface is \emph{tight} if its image meets
every open half-space in an empty or connected set.
\end{definition}

For a closed embedded surface, $\Vol(X)$ is the volume of its bounded
complementary component and $\nu_X$ is the outward unit normal.
Equal induced metrics give an isometry by the identity on the source;
images are congruent if related by a Euclidean isometry.

\begin{definition}\label{def:bending}
A field $Y:M\to\R^3$ along $X$ is an \emph{infinitesimal bending} if
\begin{equation}\label{eq:strain}
 \langle dX_p(v),dY_p(w)\rangle+
 \langle dX_p(w),dY_p(v)\rangle=0
 \quad(p\in M,\ v,w\in T_pM).
\end{equation}
Its outward normal flux is $\int_M\langle Y,\nu_X\rangle\,dA_X$.
\end{definition}

\begin{definition}\label{def:asymptotic}
A regular curve $c$ is \emph{asymptotic} if
$\mathrm{II}(c',c')=0$, equivalently if $(X\circ c)''$ is tangent
to the surface. This is independent of the regular parameter.
\end{definition}

Supports are closures in the source; $A\Subset U$ means that
$\overline A$ is compact and contained in the open set $U$.
Differences of embeddings are taken on the same source.
All $C^r$ norms use a fixed smooth background metric on that source.

\section{Classical reductions and the volume argument}
\label{sec:reductions}

\begin{theorem}[Classical curvature criterion for tightness]
\label{thm:tightness}
For a smooth embedding $X:M\to\R^3$ of a closed surface, the two-piece
property of Definition~\ref{def:tight} is equivalent to
\begin{equation}\label{eq:tightness}
 \int_M\max(K_X,0)\,dA_X=4\pi.
\end{equation}
\end{theorem}

We use this classical theorem as an external result; see
Banchoff--K\"uhnel \cite[Definition~1.1.1, p.~57]{BK}.
By Gauss--Bonnet, its equivalent
total-absolute-curvature formulation is
$\int_M|K_X|\,dA_X=8\pi-2\pi\chi(M)$.
This criterion is useful here because it reduces preservation of
tightness to preservation of the positive-curvature part. It allows
zero-curvature regions.

\begin{lemma}[Preservation under a negative-curvature modification]
\label{lem:tight-preservation}
Suppose $X_0:M\to\R^3$ is a tight embedding of a closed surface,
$U\subset\{K_{X_0}<0\}$ is open, and $X:M\to\R^3$ is an embedding
with $\supp(X-X_0)\Subset U$ and $K_X<0$ on $U$.
Then $X$ is tight.
In particular these conditions hold for every sufficiently small
$C^2$ perturbation supported in a fixed compact subset of
$\{K_{X_0}<0\}$.
\end{lemma}
\begin{proof}
On $U$ both positive curvature densities vanish. Outside the compact
support of $X-X_0$, the maps agree on a neighborhood and hence have
the same metric and curvature. Thus their total positive curvatures
are equal, and Theorem~\ref{thm:tightness} applies.
For the last assertion choose a compact neighborhood of the prescribed
support contained in $\{K_{X_0}<0\}$. On this neighborhood the original
curvature has a strictly negative upper bound. Formula
\eqref{eq:gauss-curvature} and continuity in the first two derivatives
preserve that bound under a sufficiently small $C^2$ change.
Embeddings of a compact source form an open set in the $C^1$
topology \cite[Chapter~2, \S1, Theorem~1.4]{Hirsch}.
\end{proof}

\begin{lemma}[First variation of enclosed volume]\label{lem:volume}
For an embedding $X:M\to\R^3$ of a closed surface and a smooth field
$Y:M\to\R^3$, the maps $X+tY$ are embeddings for sufficiently small
$|t|$, and
\begin{equation}\label{eq:volume-variation}
 \left.\frac{d}{dt}\right|_{t=0}\Vol(X+tY)
 =\int_M\langle Y,\nu_X\rangle\,dA_X.
\end{equation}
For small $t$, the enclosed volume is a polynomial in $t$ of degree
at most three.
\end{lemma}
\begin{proof}
The assertion about embeddings follows from the openness just cited.
Orient $M$ so that $X$ has outward normal. On $\R^3$ define the
two-form
\[
 \beta=\tfrac13(x\,dy\wedge dz+y\,dz\wedge dx+z\,dx\wedge dy),
 \qquad d\beta=dx\wedge dy\wedge dz.
\]
Set $X_t=X+tY$ and $\mathcal V(t)=\int_M X_t^*\beta$.
The divergence theorem gives $\mathcal V(0)=\Vol(X)>0$.
For every embedded $X_t$ it also gives
$\mathcal V(t)=\Vol(X_t)$ or $-\Vol(X_t)$, according to orientation.
Continuity and positivity at zero select the positive sign for
all sufficiently small $|t|$.
In positively oriented local coordinates,
\[
 X_t^*\beta=\tfrac13
       \det(X+tY,X_u+tY_u,X_v+tY_v)\,du\wedge dv.
\]
The integrand is a polynomial of degree at most three, proving
the polynomial assertion globally.

Extend the field $Y$ off $X(M)$ by a tubular projection and use
Cartan's formula for the derivative of pullback at zero.
The exact term has integral zero by Stokes' theorem on the
closed source. The remaining term is
\[
 \mathcal V'(0)=\int_M X^*(\iota_Y d\beta)
       =\int_M\langle Y,\nu_X\rangle\,dA_X.
\]
Here $\iota_Y$ denotes contraction with the extended velocity field;
its restriction to the surface is independent of the extension.
This proves \eqref{eq:volume-variation}. The classical ingredients
are Cartan's formula, Stokes' theorem, and the divergence theorem;
see \cite[Theorem~14.35 and Chapter~16, including Theorem~16.32]{Lee}.
\end{proof}

\begin{proof}[Proof of Theorem~\ref{prop:pair}]
For arbitrary tangent vectors $v,w$, expand
\[
 \langle d(X+tY)(v),d(X+tY)(w)\rangle.
\]
The linear term vanishes by \eqref{eq:strain}, leaving
$g_X(v,w)+t^2\langle dY(v),dY(w)\rangle$. This proves the exact
identity \eqref{eq:exact-metric} at $t=\pm a$.
By Lemma~\ref{lem:volume},
\begin{equation}\label{eq:odd-volume}
 \Vol(X+aY)-\Vol(X-aY)
 =2a\int_M\langle Y,\nu_X\rangle\,dA_X+O(a^3).
\end{equation}
The coefficient of $a$ is nonzero; hence the difference is nonzero
for every sufficiently small $a>0$.

Choose a compact neighborhood of $\supp Y$ inside $\{K_X<0\}$.
For fixed $Y$, $X+tY\to X$ in $C^2$, so its curvature stays negative
on that neighborhood for small $|t|$. Outside it the embedding is
unchanged. Lemma~\ref{lem:tight-preservation} proves tightness.
Euclidean isometries preserve enclosed volume, so unequal volumes
exclude congruence of the images.
\end{proof}

\section{The BGQR annulus and completion to a tight torus}
\label{sec:seed}

We prove Theorem~\ref{thm:seed} by using the annulus of
Banks--Ghomi--Qiu--Raffaelli, with a flatter choice of their end profile.

\begin{proof}[Proof of Theorem~\ref{thm:seed}]
Use the first closed asymptotic curve $\Gamma_1$ in \cite{BGQR}.
In their Section~2.4, replace the height profile on $[-200,200]$ by
\[
 H(r)=\int_0^r q(t)\,dt,
\]
where $q:\R\to[0,1]$ is smooth and even, equals one on $[-175,175]$,
is positive on $(-200,200)$, vanishes for $|r|\ge200$, and is
nonincreasing for $r\ge0$. Such a cutoff is obtained from the
normalized tail integral of a positive bump on $(175,200)$,
reflected evenly.
Then $H'>0$ in the interior, $rH''\le0$, and all positive-order
endpoint derivatives of $H$ vanish. The proofs of
\cite[Propositions~3.2--3.3]{BGQR} apply with this profile:
embeddedness uses the unchanged horizontal sections and strictly
increasing height, while their equation~(13) and end estimates~(14)
preserve $K<0$. Boundary immersion and strict convexity use only
the unchanged horizontal data. Thus we obtain a smooth embedded
annulus $S$ in $\R^2\times[-h,h]$, where $h=H(200)>0$, with
horizontal boundary tangent planes and strictly convex boundary
curves. Its interior lies in $|z|<h$. By the first paragraph of
\cite[Section~4]{BGQR}, $\Gamma_1$ lies in $|r|<175$; hence its
neighborhood there is unchanged.

Let $D_\pm$ be the planar disks bounded by $\partial S$.
Horizontal projection is a diffeomorphism on sufficiently small
boundary collars, since the boundary tangent planes are horizontal.
Each projected collar lies either inside or outside its disk;
the corresponding boundary integral of geodesic curvature, measured
against the inward conormal, is $2\pi$ or $-2\pi$.
Gauss--Bonnet gives their sum as $-\int_S K\,dA>0$ because
$\chi(S)=0$. Both collars therefore lie inside their disks.
In these graph coordinates, flat endpoint derivatives of $H$
give infinite-order contact with the boundary planes.

Choose $R$ larger than the horizontal extent of $S$ and $D_\pm$.
To provide the outer surface explicitly, choose a smooth convex
function $\phi:\R\to[0,\infty)$, zero on $(-\infty,0]$, with
$\phi',\phi''>0$ on $(0,\infty)$ and $\phi(t)\to\infty$.
For instance take the zero extension of
$\int_0^t(t-u)e^{-1/u}\,du$. The body
\[
 \Omega=\{(x,z):\phi(|x|^2-R^2)+z^2/h^2\le1\}
\]
is compact and convex, with nonempty interior and nonvanishing
defining gradient on its boundary. The boundary is a smooth
convex sphere and contains the planar caps $|x|\le R$, $z=\pm h$.
Outside these caps the defining Hessian is positive definite:
its horizontal block is $2\phi'I+4\phi''x\otimes x$ and its
vertical entry is $2/h^2$. The interior of $S$ lies strictly
inside $\Omega$, since its horizontal term is zero and $|z|<h$.
Remove the interiors of $D_\pm$ from $\partial\Omega$ and
attach $S$. At each seam the two pieces form the graph of the
annular height function on the inside and zero on the outside;
infinite-order contact makes it smooth. Strict interiority excludes
other intersections. The resulting closed embedded surface is
orientable and has Euler characteristic zero, being the union of
two annuli, hence is a torus. The removed disks were flat,
so Gauss--Bonnet on the convex sphere leaves total positive curvature
$4\pi$; $S$ adds none. Theorem~\ref{thm:tightness} proves tightness.
The retained curve remains embedded and asymptotic in $K<0$.
A parametrization by $T^2$ gives $X_0$ and $c$ as required.
\end{proof}

\section{A geometric annulus implies a supported bending}
\label{sec:bending}

We first construct the bending on a negatively curved annulus
foliated by closed asymptotic curves and straight transverse curves.


Straightness of the transverse curves supplies the vanishing
connection coefficient used below.

\subsection{The linearized metric equation in asymptotic coordinates}

For the remainder of this section use the metric $g=g_B$ and set
\[
 \mathcal W=|B_s\times B_c|=\sqrt{\det(g_{ij})}>0,\qquad
 n=\frac{B_s\times B_c}{\mathcal W}.
\]
Here $s$ is the periodic coordinate and $c$ is the transverse coordinate.
Derivatives in $s$ mean derivatives of periodic lifts; $ds$ is the
global one-form on $\R/\ell\mathbb Z$ induced by the standard one-form
on $\R$. Coordinate indices in this section range over $\{s,c\}$.
The Gauss formula and the Mainardi--Codazzi equations in these
conventions are the classical identities
\[
 B_{ij}=\Gamma^k_{ij}B_k+\mathrm{II}_{ij}n,\qquad
 (\nabla_i\mathrm{II})_{jk}=(\nabla_j\mathrm{II})_{ik};
\]
see \cite[\S4-3]{doCarmo}.

The hypotheses imply
\begin{equation}\label{eq:band-geometry}
 \mathrm{II}_{ss}=\mathrm{II}_{cc}=0,\qquad
 \sigma:=\mathrm{II}_{sc}\ne0,\qquad \Gamma^s_{cc}=0.
\end{equation}
Indeed the first diagonal entry vanishes by asymptoticity.
Since $B_c\ne0$, a transverse curve is straight precisely when
$B_{cc}$ is parallel to $B_c$; comparing with the Gauss formula
gives the other diagonal entry and the last equality.
The Gauss equation gives $K_B=-\sigma^2/\mathcal W^2$,
so $\sigma$ is nowhere zero.

Decompose a field along $B$ uniquely as
\[
 Y=dB(U)+hn,\qquad \beta=g(U,\cdot),
\]
where $U$ is a tangent vector field on the parameter annulus and
$h=Y\cdot n$. The metric dual $\beta^\sharp$ means the unique vector
field $U$ with $g(U,v)=\beta(v)$ for every tangent vector $v$.
Differentiating $\beta_j=\langle Y,B_j\rangle$ and using the
Gauss formula gives
\[
 \langle Y_i,B_j\rangle=(\nabla_i\beta)_j-h\,\mathrm{II}_{ij}.
\]
Thus the infinitesimal bending equation is exactly
\begin{equation}\label{eq:linearized-metric}
 (\nabla_i\beta)_j+(\nabla_j\beta)_i
                          -2h\,\mathrm{II}_{ij}=0.
\end{equation}
This is the classical tangential--normal formulation of the
linearized metric equation. For background on infinitesimal
bendings see \cite[Chapter~12]{Spivak}; the calculation above fixes
the formula and sign used here.

We choose $\beta=p(s,c)\,ds$. This sets the transverse covector
component to zero, leaving one scalar unknown and the normal
displacement. By \eqref{eq:band-geometry}, the three equations
in \eqref{eq:linearized-metric} become
\begin{equation}\label{eq:reduced-strain}
 p_s=\Gamma^s_{ss}p,\qquad
 0=\Gamma^s_{cc}p,\qquad
 h=\frac{p_c-2\Gamma^s_{sc}p}{2\sigma}.
\end{equation}
The middle equation is already satisfied by straightness.
The mixed equation determines $h$ algebraically.
Only the first equation remains to be solved around each circle.

\begin{proposition}[The bending family and its volume functional]
\label{prop:codazzi-family}
Under the hypotheses of Theorem~\ref{thm:band-bending}, put
\begin{equation}\label{eq:geometric-factor}
 Q=\sqrt{|\sigma|\mathcal W}=(-K_B)^{1/4}\mathcal W>0.
\end{equation}
For every $F\in C_c^\infty(I)$, define
\begin{equation}\label{eq:geometric-field}
 p=QF(c),\qquad \beta=p\,ds,\qquad
 h=\frac{p_c-2\Gamma^s_{sc}p}{2\sigma},\qquad
 Y_F=dB(\beta^\sharp)+hn.
\end{equation}
Then $Y_F$ is a smooth infinitesimal bending with
$\supp Y_F\subset(\R/\ell\mathbb Z)\times\supp F$.
Its signed normal flux is
\begin{equation}\label{eq:geometric-flux}
 \int\det(Y_F,B_s,B_c)\,ds\,dc
       =-\int_I F(c)J(c)\,dc,\qquad
 J(c)=\int_0^\ell Q\,\partial_c(\mathcal W/\sigma)\,ds.
\end{equation}
In this family a member with nonzero flux exists if and only if
$J$ is not identically zero on $I$.
\end{proposition}
\begin{proof}
The first Codazzi equation and the area-density identity give
\[
 \partial_s\log|\sigma|=\Gamma^s_{ss}-\Gamma^c_{sc},
 \qquad
 \partial_s\log\mathcal W=\Gamma^s_{ss}+\Gamma^c_{sc}.
\]
For example, the Codazzi identity
$(\nabla_s\mathrm{II})_{cs}=(\nabla_c\mathrm{II})_{ss}$,
after using the zero diagonal coefficients, reads
\[
 \sigma_s-(\Gamma^c_{sc}+\Gamma^s_{ss})\sigma
                                      =-2\Gamma^c_{sc}\sigma.
\]
The area identity follows by differentiating
$\mathcal W=\sqrt{\det g}$ and using the metric connection.
Adding the logarithmic identities yields
\[
 Q_s/Q=\Gamma^s_{ss}.
\]
Consequently $p=QF(c)$ solves the first equation in
\eqref{eq:reduced-strain}. The other two equations hold by
construction, so \eqref{eq:linearized-metric} proves zero strain.
The factor $Q$ is smooth, positive, and periodic.

All coefficients are smooth and periodic, and $h$ involves only
$F$ and $F'$, proving the support assertion. The cutoff belongs in
the free profile $F$; multiplying a bending by an arbitrary cutoff
would introduce extra strain terms.

For the flux, the second Codazzi equation gives
\[
 \partial_c\log|\sigma|=\Gamma^c_{cc}-\Gamma^s_{sc},\qquad
 \partial_c\log\mathcal W=\Gamma^s_{sc}+\Gamma^c_{cc}.
\]
Its first equality follows by expanding
$(\nabla_c\mathrm{II})_{sc}=(\nabla_s\mathrm{II})_{cc}$.
Subtracting gives
\[
 \partial_c(\mathcal W/\sigma)
             =2\Gamma^s_{sc}(\mathcal W/\sigma).
\]
The normal displacement in \eqref{eq:geometric-field} therefore
satisfies
\[
 h\mathcal W=\tfrac12\left\{
       \frac{\mathcal W}{\sigma}p_c
       -\partial_c\left(\frac{\mathcal W}{\sigma}\right)p
                         \right\}.
\]
Integrating by parts in $c$ removes the total derivative, because
$p$ has compact support in that direction. We obtain
\[
 \int h\,dA_B
       =-\int p\,\partial_c(\mathcal W/\sigma)\,ds\,dc,
\]
which is \eqref{eq:geometric-flux}.

If $J\equiv0$, every member of this family has zero flux.
If $J(c_*)\ne0$, choose a nonnegative smooth cutoff $\chi$
supported in a neighborhood of $c_*$ and positive at $c_*$.
The profile $F=\chi J$ then gives flux
$-\int_I\chi J^2\,dc<0$. This proves the final equivalence.
\end{proof}

Nonzero normal displacement need not give nonzero flux.
The second Codazzi identity also gives
\[
 h(-K_B)\mathcal W
      =\tfrac12\partial_c(p\sigma/\mathcal W).
\]
Its integral along each ruling vanishes by compact support.
We next show that the unweighted flux is nonzero for some profile.

\subsection{Nonzero flux}

\begin{proof}[Proof of Theorem~\ref{thm:band-bending}]
Fix the prescribed open subinterval $I_0$ and choose a reference
leaf inside it. Choose $s=0$ as one fixed ruling. Relabel nearby
leaves by signed Euclidean distance along that straight ruling,
with the reference leaf labeled $c=0$ and the positive direction
agreeing with the original transverse orientation.
Restrict to a neighborhood $|c|<\varepsilon$ whose closure is
inside the chosen subband, for some $\varepsilon>0$.
This change of leaf label is smooth, increasing, and independent
of $s$, because the ruling is regular and straight.
It preserves the hypotheses, compact supports, and signed flux
under change of variables. We retain the notation $B(s,c)$;
all quantities below, including $J$, are computed in these labels.

Put $\gamma(s)=B(s,0)$ and $V(s)=B_c(s,0)$.
Straightness gives a unique smooth affine ruling coefficient $u$:
\[
 B(s,c)=\gamma(s)+u(s,c)V(s),\qquad
 u(s,0)=0,\quad u_c(s,0)=1.
\]
For example $u=\langle B-\gamma,V\rangle/|V|^2$.
After decreasing $\varepsilon$, compactness of the $s$-circle
gives $u_c>0$ throughout the band.
The normalization on the selected ruling gives
$|V(0)|=1$ and $u(0,c)=c$.
The distance normalization applies only at $s=0$.

Set $\Delta(s)=\det(\gamma_s,V,V_s)$.
At $c=0$ the mixed second coefficient has unnormalized value
$\Delta$, and the two diagonal coefficients vanish. Hence
\[
 K_B(s,0)=-\frac{\Delta(s)^2}{|\gamma_s\times V|^4}<0,
\]
so $\Delta$ is nowhere zero. Its sign is constant on the circle.

Expanding the asymptotic condition
$\det(B_{ss},B_s,V)=0$ gives
\[
 \det(\gamma_{ss}+uV_{ss},\,\gamma_s+uV_s,\,V)
                                  +2\Delta u_s=0.
\]
The first determinant is a polynomial of degree at most two in $u$.
Its constant coefficient is zero because $\gamma$ is asymptotic.
Differentiate the displayed equality in $c$ at $c=0$.
Since $u_c(s,0)=1$ and $u_{sc}(s,0)=0$, its linear
coefficient is also zero.
Consequently, on $c>0$, $u_s/u^2$ depends only on $s$.
Commutation of mixed derivatives now gives
\[
 \partial_s(u_c/u^2)=\partial_c(u_s/u^2)=0.
\]
Evaluate this $s$-independent quantity at $s=0$, where
$u(0,c)=c$ and $u_c(0,c)=1$. We obtain the geometric
normalization identity
\begin{equation}\label{eq:ruling-normalization}
 c^2u_c=u^2\qquad(0<c<\varepsilon).
\end{equation}

Use the squared area vector in the affine ruling parameter:
\[
 S(s,u)=|(\gamma_s+uV_s)\times V|^2,\qquad
 S_u(s,u)=S_1(s)+2S_2(s)u,\qquad
 S_2=|V_s\times V|^2>0.
\]
Here $S$ is quadratic in $u$, and $S_1(s)=S_u(s,0)$;
the strict positivity follows from $\Delta\ne0$.
From $B_c=u_cV$ and $B_{sc}=u_{cs}V+u_cV_s$,
\[
 \mathcal W^2=u_c^2S,\qquad
 \sigma=\frac{u_c^2\Delta}{\mathcal W},\qquad
 \frac{\mathcal W}{\sigma}=\frac S\Delta,\qquad
 Q=\sqrt{|\Delta|}\,u_c .
\]
Thus the flux kernel from \eqref{eq:geometric-flux}
satisfies, for $0<c<\varepsilon$,
\begin{equation}\label{eq:direct-flux}
 \begin{aligned}
 J(c)&=\operatorname{sign}(\Delta)
       \int_0^\ell\frac{u_c^2S_u}{\sqrt{|\Delta|}}\,ds\\
 &=\frac{\operatorname{sign}(\Delta)}{c^4}
       \int_0^\ell
       \frac{u^4}{\sqrt{|\Delta|}}(S_1+2S_2u)\,ds .
 \end{aligned}
\end{equation}

Suppose $J$ vanished identically on this positive subinterval.
The last integral would then be zero for every such $c$.
Repeatedly apply $c^2\partial_c$ and use
\eqref{eq:ruling-normalization}. We obtain by induction
\begin{equation}\label{eq:vanishing-moments}
 \int_0^\ell\frac{u^{n+4}}{\sqrt{|\Delta|}}
       \left(S_1+\frac{n+4}{2}S_2u\right)\,ds=0,
 \qquad n=0,1,2,\ldots .
\end{equation}
For clarity, applying $c^2\partial_c$ to the expression with
index $n$ gives $n+4$ times the expression with index $n+1$.
The coefficient functions depend only on $s$.
Smoothness and compactness of the circle
justify differentiation under the integral.

Fix any one leaf with $0<c<\varepsilon$.
On its compact circle, $u>0$ and $S_2>0$, while $S_1$ is
bounded. Choose an integer $n$ large enough that
\[
 (n+4)\min_s(S_2u)>2\max_s|S_1|.
\]
The parenthesis in \eqref{eq:vanishing-moments} is then strictly
positive at every $s$, and so is the integral, a contradiction.

It follows that $J$ is not identically zero.
Choose a nonnegative smooth cutoff $\chi$, compactly supported
inside $(0,\varepsilon)$ and positive at a point where $J\ne0$,
and set $F=\chi J$ there, extended by zero.
Proposition~\ref{prop:codazzi-family} gives a bending with flux
$-\int\chi J^2\,dc<0$ and compact support in the selected subband.
Return to the original leaf labels and extend the field by zero
over the rest of the original annulus. Compact support makes
this extension smooth; the chain rule preserves zero strain,
and the orientation-preserving change of variables preserves
the signed flux. Its support lies compactly inside the original
$(\R/\ell\mathbb Z)\times I_0$, as required.
\end{proof}

\section{The Darboux frame and the forced inflection}\label{sec:frame}

We now construct the annulus needed for
Corollary~\ref{thm:one-curve}, starting with its given asymptotic curve.

Let $\gamma:\R/\ell\mathbb Z\to\R^3$, $\ell>0$, be the given curve
parametrized by arclength, and let $n(s)$ be the unit normal of
the surface along it. Put $T=\gamma_s$ and $E=n\times T$.
Then $(T,E,n)$ is a positively oriented orthonormal frame.
All derivatives in $s$ are derivatives of periodic lifts.

\begin{lemma}[Frame equations]\label{lem:frame}
There are smooth real periodic functions $k,\tau$ satisfying
\begin{equation}\label{eq:darboux}
 T_s=kE,\qquad E_s=-kT+\tau n,\qquad n_s=-\tau E.
\end{equation}
They are given by $k=\langle T_s,E\rangle$ and
$\tau=\langle E_s,n\rangle$. Moreover $K=-\tau^2$ along $\gamma$.
\end{lemma}
\begin{proof}
Unit speed gives $T_s\perp T$, and asymptoticity gives
$T_s\perp n$, proving the first equation.
Differentiating orthonormality gives the other two, with the
$T$ component of $n_s$ equal to
$-\langle n,T_s\rangle=0$.
In the orthonormal tangent basis $(T,E)$, the second fundamental
form has first entry zero and mixed entry $\tau$.
Its determinant is therefore $-\tau^2$, which proves the curvature
identity by \eqref{eq:gauss-curvature}.
\end{proof}

Here $k$ is the signed geodesic curvature and $\tau$ is the signed
geodesic torsion. Negative Gaussian curvature implies $\tau\ne0$;
hence its sign is constant on the circle. Define
\begin{equation}\label{eq:rho}
 \rho=|\tau|^{-1/2}>0.
\end{equation}
This smooth periodic function equals the negative Gaussian curvature
along the core raised to the power $-1/4$.
We use this surface frame even
at points where the curve's curvature vanishes, where a Frenet
normal need not be defined.
The exponent is chosen so that
$\rho_s/\rho=-\tau_s/(2\tau)$, the integrating factor relation
in the asymptotic return equation of Section~\ref{sec:ruled}.

\begin{definition}\label{def:inflection}
An \emph{inflection} of the unit-speed curve $\gamma$ is a point
$s_0$ where $\gamma_{ss}(s_0)=0$. By \eqref{eq:darboux} this is
equivalent to $k(s_0)=0$. No prescribed order of vanishing is
included in this definition.
\end{definition}

\begin{lemma}[Injectivity of the negative Gauss map]\label{lem:gauss}
For a smooth embedded tight torus with outward normal $n$, the
restriction of $n$ to $\{K<0\}$ is injective.
\end{lemma}
\begin{proof}
For a regular value $u\in S^2$ of $n$, let $m_+(u)$ and $m_-(u)$
be the numbers of preimages with positive and negative curvature.
They are finite because the source is compact and each preimage
is isolated. A maximum of the height function $\langle X,u\rangle$
has outward normal $u$. At that point its Hessian is negative
semidefinite; regularity of $u$ makes the Hessian nondegenerate,
so it is negative definite and the Gaussian curvature is positive.
Thus $m_+(u)\ge1$ for every regular value.

The change-of-variables formula for the Gauss map, whose absolute
Jacobian is $|K|$, and Theorem~\ref{thm:tightness} give
\[
 \int_{S^2}m_+(u)\,dA(u)=\int_{\{K>0\}}K\,dA=4\pi .
\]
Sard's theorem makes the set of nonregular values have measure zero.
Since $S^2$ has area $4\pi$, it follows that $m_+=1$ almost
everywhere. The signed-degree formula gives
\[
 m_+(u)-m_-(u)=\deg n
 =\frac1{4\pi}\int_{T^2}K\,dA=\frac{\chi(T^2)}2=0.
\]
Here the equality for the degree follows either from
$n^*(dA_{S^2})=K\,dA$ or the usual integration formula for degree,
and the next equality is Gauss--Bonnet. Thus $m_-=1$ almost
everywhere as well. Change of variables, Sard's theorem, and the
degree formula are used in their ordinary smooth-manifold forms;
see \cite[Chapters~6 and~17]{Lee}.

If two different points in $\{K<0\}$ had the same normal, the inverse
function theorem would provide disjoint neighborhoods mapped
diffeomorphically onto overlapping open subsets of $S^2$.
Every regular value in their nonempty open overlap would have at
least two negative-curvature preimages. This contradicts $m_-=1$
almost everywhere, since a nonempty open subset of $S^2$ has
positive area.
\end{proof}

\begin{proposition}[One loop supplies an inflection]\label{prop:inflection}
Every regular embedded closed asymptotic curve in $\{K<0\}$ on
a smooth embedded tight torus has an inflection. In fact its
signed geodesic curvature $k$ takes both positive and negative values.
\end{proposition}
\begin{proof}
Lemma~\ref{lem:gauss} shows that $s\mapsto n(s)$ is injective.
It is regular because $|n_s|=|\tau|>0$. It is therefore a simple
regular closed curve on the unit sphere.

We use the Gauss--Bonnet--Fenchel argument recorded by
Ghomi--Raffaelli \cite[Note~5.2, equations~(10)--(12)]{GR}, where
it is attributed to Kovaleva \cite{Kovaleva}. We give the calculation
to specify its orientations and its application here.
Let $\sigma$ be arclength on the spherical normal curve, and
let $\kappa_{\mathrm{sph}}$ be its signed geodesic curvature
for the outward orientation of $S^2$.
The frame equations give
\[
 n_{ss}=-\tau_sE+\tau kT-\tau^2n,\qquad
 n\times n_s=\tau T.
\]
Consequently
\[
 \kappa_{\mathrm{sph}}
 =\frac{\langle n_{ss},n\times n_s\rangle}{|n_s|^3}
 =\frac{k}{|\tau|},\qquad d\sigma=|\tau|\,ds,
\]
and therefore $\int\kappa_{\mathrm{sph}}\,d\sigma=\int_0^\ell k\,ds$.

Choose the disk on the sphere lying to the left of the oriented
normal curve, and let its area be $A$. The Jordan curve theorem
on the sphere gives two disks with nonempty open interiors, so
$0<A<4\pi$. The sphere has curvature one and this disk has Euler
characteristic one. Gauss--Bonnet with boundary gives
\[
 \int\kappa_{\mathrm{sph}}\,d\sigma+A=2\pi,
 \qquad
 \left|\int_0^\ell k\,ds\right|=|2\pi-A|<2\pi .
\]
Gauss--Bonnet with boundary is used in its standard smooth-boundary
form \cite[\S4-5]{doCarmo}. No convexity or hemisphere condition
on the spherical disk is required.

On the other hand Fenchel's theorem \cite{Fenchel} for a regular
closed space curve gives
\[
 \int_0^\ell|\gamma_{ss}|\,ds
 =\int_0^\ell|k|\,ds\ge2\pi.
\]
If $k$ had only one sign, including the possibility of zeros,
the last integral would equal $|\int k\,ds|$, a contradiction.
Thus $k$ has both signs, and continuity gives a zero.
\end{proof}

The role of the inflection is to permit a correction concentrated
near a point where $k=0$ without significantly changing the
transverse component of a ruling. Before choosing that ruling,
we arrange one further derivative condition. The frame
equations and $\tau_s/\tau=-2\rho_s/\rho$ give
\[
 (\rho n)_{ss}=(\rho_{ss}-\rho^{-3})n+\rho\tau kT.
\]
Thus at an inflection the condition
$\rho_{ss}\ne\rho^{-3}$ means exactly that the scaled normal
curve $\rho n$ has nonzero acceleration. Here the scaled normal
is the surface normal multiplied by the curvature weight
$(-K)^{-1/4}$ along the core. The next section arranges the
nonzero-acceleration condition by a local modification.

\section{Local preparation of the curve}\label{sec:preparation}

\begin{proposition}[An inflection implies a prepared curve]
\label{prop:prepare}
Let $\ell>0$, let $X_0:M\to\R^3$ be a smooth embedding of a closed
oriented surface, and let $c:\R/\ell\mathbb Z\to M$ be a smooth embedding
such that $\gamma=X_0\circ c$ is a unit-speed asymptotic curve
in $\{K_{X_0}<0\}$. Use its frame from
Section~\ref{sec:frame}, and suppose $k(s_0)=0$.
For every open neighborhood $U$ of $c(\R/\ell\mathbb Z)$ with
$U\subset\{K_{X_0}<0\}$ and every $\varepsilon>0$, there is a
smooth embedding $\widetilde X:M\to\R^3$ such that
\[
 \supp(\widetilde X-X_0)\Subset U,\quad
 \|\widetilde X-X_0\|_{C^2}<\varepsilon,\quad
 K_{\widetilde X}<0\text{ on }U.
\]
The curve $\widetilde X\circ c$ is closed, embedded, and
asymptotic. In its own arclength $\widetilde s$, chosen to be zero at
$c(s_0)$, its frame functions satisfy
\begin{equation}\label{eq:prepared}
 \widetilde k(0)=0,\qquad
 \widetilde\rho_{\widetilde s\widetilde s}(0)\ne\widetilde\rho(0)^{-3}.
\end{equation}
Here $\widetilde\rho=|\widetilde\tau|^{-1/2}$ for the new curve.
If $X_0$ is tight, $\widetilde X$ is tight.
\end{proposition}
\begin{proof}
We change the curve while retaining its tangent directions
and its normal directions. Prescribe a new velocity $v(s)T(s)$
with $v>0$. To keep the curve closed it is necessary and sufficient
that
\begin{equation}\label{eq:closure-moment}
 \int_0^\ell(v-1)T\,ds=0,
\end{equation}
because $\int_0^\ell T\,ds=0$. We construct a family $v=v_\delta$,
$\delta\downarrow0$, with
\begin{equation}\label{eq:speed-jets}
 v(s_0)=1,\quad v_s(s_0)=0,\quad v_{ss}(s_0)=d,
 \qquad \|v-1\|_{C^1}\longrightarrow0,
\end{equation}
where the fixed real number $d$ will be chosen below.

There is an interval disjoint from $s_0$ on which $k\ne0$.
Indeed $k\equiv0$ would make $T$ constant, contradicting
$\int T\,ds=0$ and $|T|=1$; a nonzero value has an open
neighborhood away from the zero at $s_0$. On such an interval,
\[
 T_s=kE,\quad T_{ss}=-k^2T+k_sE+k\tau n,\quad
 \det(T,T_s,T_{ss})=k^2\tau\ne0.
\]
It follows that the vectors $T(s)$ on that interval span $\R^3$:
otherwise they and their first two derivatives would lie in a
fixed proper linear subspace, contradicting the determinant.
Choose three points with linearly independent tangent vectors,
and choose smooth bumps $b_1,b_2,b_3$ of integral one in
disjoint sufficiently small neighborhoods of those points.
The three vectors $\int b_iT\,ds$ are then linearly independent.

Choose a smooth cutoff $\chi$ on $\R$, equal to one near zero
and supported in $(-1,1)$, and set in a coordinate interval at $s_0$
\[
 h_\delta(s)=\tfrac d2(s-s_0)^2
                        \chi((s-s_0)/\delta),
\]
extended by zero periodically. Its value and first derivative at
$s_0$ are zero, and its second derivative there is $d$.
Moreover $\|h_\delta\|_{C^1}=O(\delta)$ and
$|\int h_\delta T\,ds|=O(\delta^3)$.
Solve the invertible three-dimensional linear system
\[
 \sum_{i=1}^3 a_i\int b_iT\,ds=-\int h_\delta T\,ds.
\]
Its matrix is fixed, so $a_i=O(\delta^3)$.
Then $v=1+h_\delta+\sum_i a_i b_i$ satisfies
\eqref{eq:closure-moment} and \eqref{eq:speed-jets}.
For small $\delta$ it is positive.

Integrate $\widetilde\gamma_s=vT$ with
$\widetilde\gamma(s_0)=\gamma(s_0)$. The moment condition proves
periodicity; estimates on $v-1$ and $v_s$ show
$\widetilde\gamma-\gamma\to0$ in $C^2$.
Since the surface is oriented, its embedded source circle has a
two-sided smooth collar. Choose coordinates $(s,w)$ in such a
collar, with $X_0(s,0)=\gamma(s)$, whose closure lies in $U$.
For a fixed cutoff $\chi_0(w)$ equal to one near zero and
supported in the collar, define
\[
 \widetilde X(s,w)=X_0(s,w)+
              \chi_0(w)(\widetilde\gamma(s)-\gamma(s)),
\]
and put $\widetilde X=X_0$ elsewhere.
This tends to $X_0$ in $C^2$. Embedding and negative curvature
persist, and Lemma~\ref{lem:tight-preservation} preserves
tightness when present.

At $w=0$, the two tangent vectors of $\widetilde X$ are $vT$
and the original transverse tangent. Their span and orientation
are unchanged; hence the normal is exactly $n(s)$.
The new acceleration
$\widetilde\gamma_{ss}=v_sT+vkE$ is tangent to this plane,
so the new curve is asymptotic.
Let $\widetilde s$ be its arclength, with $\widetilde s(s_0)=0$.
Its new period is $\widetilde\ell=\int_0^\ell v(s)\,ds$.
Then
\[
 d\widetilde s=v\,ds,\qquad
 \widetilde k=k/v,\qquad
 \widetilde\tau=\tau/v,\qquad
 \widetilde\rho=\sqrt v\,\rho.
\]
At $s_0$, where $v=1$ and $v_s=0$, differentiating twice yields
\[
 \widetilde k=0,\qquad
 \widetilde\rho_{\widetilde s\widetilde s}=\rho_{ss}+\tfrac12\rho d,
 \qquad
 \widetilde\rho^{-3}-\widetilde\rho_{\widetilde s\widetilde s}
 =\rho^{-3}-\rho_{ss}-\tfrac12\rho d .
\]
Choose $d$ different from
$2(\rho^{-3}-\rho_{ss})(s_0)/\rho(s_0)$.
The last expression is then nonzero, proving
\eqref{eq:prepared} in the new arclength.
\end{proof}

\begin{proof}[Proof of Theorem~\ref{thm:prepared-core}]
Reparametrize the original core by arclength.
Proposition~\ref{prop:inflection} gives an inflection, and
Proposition~\ref{prop:prepare} gives the required tight embedding
$X_1$. Reparametrize the new curve by its own arclength to obtain
$c_1$. Along an asymptotic curve, $\rho=(-K)^{-1/4}$ and the frame
equations give
$(\rho n)_{ss}=(\rho_{ss}-\rho^{-3})n+\rho\tau kT$.
Thus the conclusion of Proposition~\ref{prop:prepare} gives
$\gamma_{ss}=0$ and $(\rho n)_{ss}\ne0$ at the selected point.
\end{proof}

\section{From a prepared curve to a periodic ruled annulus}
\label{sec:ruled}

We construct a ruled strip around the prepared curve, then insert
it into the surface in Section~\ref{sec:insertion}.

\begin{proposition}[Prepared curve implies periodic ruled annuli]
\label{prop:ruled}
Let $\ell>0$, and suppose $\gamma:\R/\ell\mathbb Z\to\R^3$ is a smooth embedded
unit-speed curve with a smooth positively oriented orthonormal frame
$(T,E,n)$ satisfying $T=\gamma_s$, \eqref{eq:darboux}, and
$\tau\ne0$. Put $\rho=|\tau|^{-1/2}$.
If for some $s_0$
\[
 k(s_0)=0,\qquad
 \rho(s_0)^{-3}-\rho_{ss}(s_0)\ne0,
\]
then there exist smooth periodic functions $f>0$, $g$, and $\omega$
such that, for $V=gT+fE$:
\begin{enumerate}
 \item $R(s,u)=\gamma(s)+uV(s)$ is an embedded negatively curved
 ruled strip on $|u|<a$ for some $a>0$, with the prescribed tangent
 plane $\operatorname{span}(T,E)$ along $u=0$;
 \item $\det(\gamma_s,V,V_s)=f^2\tau\ne0$;
 \item for all sufficiently large $c$,
 \[
 B(s,c)=\gamma(s)+\frac{\rho(s)}{c+\omega(s)}V(s)
 \]
 has closed asymptotic constant-$c$ curves.
 For every bounded nonempty open interval of such $c$, this is an embedding of
 the corresponding open annulus. Its constant-$c$ leaves
  converge uniformly to $\gamma$ as $c\to+\infty$.
\end{enumerate}
\end{proposition}

\begin{proof}
Write a transverse vector in the prescribed tangent plane as
$V=gT+fE$, where $f>0$. We choose its direction and scale
jointly. Allowing the transverse component $f$ to vary will turn
the period correction into linear equations. Impose
\begin{equation}\label{eq:ruling-constraint}
 f_s=-kg.
\end{equation}
This removes the $E$ component of $V_s$. It will also remove
the multiplicative obstruction in the return map calculated below.
For the moment $f,g$ are unspecified. Define the geometric component
$\xi=\langle V_s,T\rangle=g_s-kf$. Then
\[
 V_s=\xi T+f\tau n,\qquad
 V_{ss}=\xi_sT+(k\xi-f\tau^2)E+(f_s\tau+f\tau_s)n.
\]
In particular
\begin{equation}\label{eq:delta}
 \Delta:=\det(T,V,V_s)=f^2\tau\ne0.
\end{equation}
At $u=0$ the ruling is transverse to the core, so $R$ is an
immersion near it. Its mixed unnormalized second coefficient is
$\langle R_{su},R_s\times R_u\rangle=\Delta$ and $R_{uu}=0$.
Formula \eqref{eq:gauss-curvature} therefore gives
\begin{equation}\label{eq:ruled-curvature}
 K_R=-\frac{\Delta^2}{|R_s\times R_u|^4}<0
\end{equation}
where $R$ is immersed.

We determine the conditions for the nonruling asymptotic curves to
close. Their graphs are $s\mapsto R(s,u(s))$.
The acceleration has normal component equal to that of
$R_{ss}+2u_sR_{su}$: the $u_{ss}R_u$ term is tangent and
$R_{uu}=0$. Set
\[
 2A=\xi_s-\frac{\tau_s}{\tau}\xi+g\tau^2 .
\]
The frame equations give the two unnormalized coefficients
\[
 \langle R_{ss},R_s\times V\rangle
     =f^2\tau_su-2f^2\tau Au^2,\qquad
 \langle R_{su},R_s\times V\rangle=f^2\tau.
\]
Thus the asymptotic equation is the Riccati equation
\begin{equation}\label{eq:riccati}
 u_s=\frac{\rho_s}{\rho}u+Au^2.
\end{equation}
Here $\tau_s/\tau=-2\rho_s/\rho$ follows from
$\rho=|\tau|^{-1/2}$. For $u\ne0$, taking the reciprocal gives
\begin{equation}\label{eq:reciprocal}
 (\rho/u)_s=-\rho A.
\end{equation}
The reciprocal substitution is the standard Bernoulli reduction
\cite[\S1.4]{Teschl}. It also identifies the exact return map.
Take $s=0$ as base point and let $\mathcal P(u_0)$ be the value
at $s=\ell$ of the solution with initial value $u_0$.
For all sufficiently small $u_0$, integration of
\eqref{eq:reciprocal} gives
\begin{equation}\label{eq:return-map}
 \mathcal P(u_0)=
 \frac{u_0}{1-\dfrac{u_0}{\rho(0)}
                      \int_0^\ell\rho A\,ds}.
\end{equation}
Small initial values keep the denominator nonzero throughout
one period. The value at $u_0=0$ is defined by continuity.
The imposed condition $f_s+kg=0$ has made the linear coefficient
in \eqref{eq:riccati} equal to $\rho_s/\rho$, so its return
multiplier is one. Accordingly $\mathcal P'(0)=1$ even when
the displayed integral is nonzero.
The whole return map is the identity on an interval precisely
when $\int_0^\ell\rho A\,ds=0$, whereas a unit multiplier
alone does not imply periodicity of neighboring solutions.

It follows that periodicity of the neighboring nonzero solutions
is equivalent to that single remaining integral condition. Indeed integration over a period proves
necessity; if this integral vanishes, a periodic function $\omega$
with $\omega_s=-\rho A$ exists, and every nonzero solution with
positive denominator is
\begin{equation}\label{eq:leaves}
 u(s,c)=\frac{\rho(s)}{c+\omega(s)}.
\end{equation}
The parameter $c$ is the additive constant in $\rho/u$; adding a
constant to $\omega$ only changes this label.

It remains to choose $f,g$ so that both
\eqref{eq:ruling-constraint} and the period condition hold.
The point of retaining the scale $f$ is that these requirements
become linear. Since $\tau^2=\rho^{-4}$ and
$\tau_s/\tau=-2\rho_s/\rho$, two integrations by parts give
\begin{equation}\label{eq:period-identity}
 \begin{aligned}
 2\int_0^\ell\rho A\,ds
 &=\int_0^\ell(\rho\xi_s+2\rho_s\xi+g\rho^{-3})\,ds\\
 &=\int_0^\ell\{\rho_s(g_s-kf)+g\rho^{-3}\}\,ds\\
 &=\int_0^\ell\{g(\rho^{-3}-\rho_{ss})-fk\rho_s\}\,ds .
 \end{aligned}
\end{equation}
There are no boundary terms because every function is periodic.
Consequently it suffices to solve
\begin{equation}\label{eq:linear-periods}
 f_s+kg=0,\qquad
 \int_0^\ell\{g(\rho^{-3}-\rho_{ss})-fk\rho_s\}\,ds=0,
 \qquad f>0.
\end{equation}

The prepared inflection separates the effects of a local correction.
Changing $g$ changes $f$ through $f_s=-kg$; if the change is
concentrated near $k(s_0)=0$, its accumulated effect on $f$
tends to zero. In the second equation of
\eqref{eq:linear-periods}, the same concentration is weighted
by $\rho^{-3}-\rho_{ss}$, whose value at $s_0$ is nonzero.
Thus a finite correction of the period can coexist with $f$
remaining uniformly close to one and positive.
The following construction implements these two effects.

Since a closed unit-speed curve
cannot have constant tangent, $\int_0^\ell k^2\,ds>0$.
Let $\varphi_\delta$ be a nonnegative smooth function of integral
one supported in the $\delta$-neighborhood of $s_0$ on the circle.
Define
\[
 b_\delta=\varphi_\delta-
       \frac{\int_0^\ell k\varphi_\delta\,ds}
            {\int_0^\ell k^2\,ds}\,k,\qquad
 (P_\delta)_s=kb_\delta,\qquad P_\delta(s_*)=0,
\]
where $s_*$ is any fixed base point.
The integral of $kb_\delta$ is zero, so $P_\delta$ is periodic.
Since $k(s_0)=0$,
\[
 \int_0^\ell|kb_\delta|\,ds\longrightarrow0,\quad
 \|P_\delta\|_\infty\longrightarrow0,\quad
 \int_0^\ell b_\delta F\,ds\longrightarrow F(s_0)
\]
for every fixed continuous periodic $F$.
The first assertion follows by bounding the bump contribution by
$\sup_{|s-s_0|<\delta}|k(s)|$ and the correction by a constant
times $|\int k\varphi_\delta|$; both tend to zero.
The second follows by integration, and the third follows from
concentration of $\varphi_\delta$ and the vanishing correction.

Set
\begin{equation}\label{eq:balance}
 \lambda_\delta=
 -\frac{\int_0^\ell k\rho_s\,ds}
 {\int_0^\ell\{b_\delta(\rho^{-3}-\rho_{ss})
                         +P_\delta k\rho_s\}\,ds},
 \qquad
 f=1+\lambda_\delta P_\delta,\quad g=-\lambda_\delta b_\delta .
\end{equation}
The denominator converges to
$\rho(s_0)^{-3}-\rho_{ss}(s_0)\ne0$.
Thus $\lambda_\delta$ is bounded for small $\delta$, and
$f\to1$ uniformly, giving $f>0$. The first equality in
\eqref{eq:linear-periods} follows by differentiating $f$.
The second follows by substitution of \eqref{eq:balance}.
Fix one such $\delta$ once and for all. This fixes the smooth
ruling $V$, all its derivatives, and a periodic $\omega$.

Compactness first gives a uniform width on
which $R$ is immersed. It is a local embedding at every $(s,0)$
by transversality. If no narrower strip were injective, there would
be pairs of distinct parameters with $|u|,|u'|\to0$ and equal
images. Subsequences of their circle coordinates converge.
Equality of the limiting core points and embeddedness of $\gamma$
force both limits to be the same $s$. A local embedding near
$(s,0)$ then excludes those pairs, a contradiction.
After narrowing once more, injectivity on the compact closed
cylinder gives a homeomorphism onto its image. Its restriction
to the interior is therefore an embedded strip of some width $a>0$.

Finally, for $c$ sufficiently large, $c+\omega>0$ and
$0<u(s,c)<a$, uniformly in $s$. The map
$(s,c)\mapsto(s,u(s,c))$ is a diffeomorphism onto its image because
\[
 u_c=-\frac{\rho}{(c+\omega)^2}<0
\]
and its inverse is $c=\rho(s)/u-\omega(s)$.
Composing with $R$ proves the asserted annular embedding.
Equation \eqref{eq:reciprocal} proves that its constant-$c$ circles
are asymptotic. Their distance from the core tends uniformly to
zero as $c\to+\infty$.
\end{proof}

\begin{proof}[Proof of Theorem~\ref{thm:ruled-interface}]
Use the Darboux frame of $\gamma$. Lemma~\ref{lem:frame} gives
$|\tau|=\sqrt{-K}$, so the curvature weight is
$\rho=|\tau|^{-1/2}$. At an inflection,
$(\rho n)_{ss}=(\rho_{ss}-\rho^{-3})n$.
The hypotheses therefore imply those of Proposition~\ref{prop:ruled},
which supplies the strip. Its constant-$c$ leaves converge uniformly
to the core, so a sufficiently large bounded leaf interval gives
the asserted annulus in any prescribed neighborhood $N$.
\end{proof}

\section{Insertion and passage to a global bending}
\label{sec:insertion}

The ruled strip just constructed has the correct core and tangent
planes, but need not already be part of the original surface.
We state the replacement step for an arbitrary embedded strip
with those first-order data. Straightness is needed in the
construction of the bending annulus, but not in this replacement
argument. The curvature estimate is needed because cutting off
a quadratic normal displacement need not be small in $C^2$.

\begin{proof}[Proof of Theorem~\ref{prop:insert}]
Use an ambient tubular neighborhood of $X_0(M)$ with its normal
projection $\pi$ onto $X_0(M)$; see
\cite[Theorem~6.24]{Lee}. After narrowing the strip, define
\[
 q=X_0^{-1}\circ\pi\circ Z .
\]
At $u=0$, $q(s,0)=c(s)$, and $dq$ is invertible because
$Z_s,Z_u$ span the original tangent plane.
Thus $q$ is a local diffeomorphism near the core. It is a
diffeomorphism onto a source neighborhood for a sufficiently
narrow uniform strip. For uniform injectivity, any pairs of equal images with strip
widths tending to zero would have limiting points on the same
embedded core. Those limits must coincide, and the inverse
function theorem at that point excludes the pairs.
After a further restriction to a smaller closed strip, compactness
also gives the required homeomorphism onto the image.
By the definition of normal projection, in that neighborhood
\[
 Z\circ q^{-1}=X_0+hn
\]
for a smooth scalar function $h$ and the original surface normal
$n$. Equality of points and tangent planes along the core gives
$h=dh=0$ there.

Fix an original-source collar $(s,w)$ of the core with compact
closure in $U$. Narrow it to lie in the image of $q$. Here $s$
agrees with the given core parameter on $w=0$.
Taylor's formula with integral remainder gives
\[
 h(s,w)=w^2b_0(s,w)
\]
with $b_0$ smooth. Take a fixed smooth cutoff $\chi$ equal to one
on $[-1,1]$ and zero outside $(-2,2)$. For small $\eta>0$ put
\[
 h_\eta(s,w)=\chi(w/\eta)w^2b_0(s,w),\qquad
 X_\eta=X_0+h_\eta n
\]
on the collar and $X_\eta=X_0$ elsewhere.
For $\eta$ small the support lies compactly in $U$.
All the following estimates are uniform on the collar, with constants
depending on the already fixed strip:
\[
 h_\eta=O(\eta^2),\quad
 (h_\eta)_s=O(\eta^2),\quad(h_\eta)_w=O(\eta),
\]
\[
 (h_\eta)_{ss}=O(\eta^2),\quad
 (h_\eta)_{sw}=O(\eta),\quad(h_\eta)_{ww}=O(1).
\]
Thus $X_\eta\to X_0$ in $C^1$, giving immersion and global
embedding for small $\eta$. Its consistently chosen unit normal
satisfies $n_\eta-n=O(\eta)$.

Let $L_\eta,M_\eta,N_\eta$ be the second fundamental coefficients
of $X_\eta$ in the fixed collar. For the original surface
$L_0(s,0)=0$ by asymptoticity, so $L_0(s,w)=O(w)$.
Also $M_0(s,0)$ is nowhere zero, because at the core the
determinant of the second fundamental form is
$-M_0(s,0)^2<0$. The derivative estimates just stated and
$n_\eta-n=O(\eta)$ imply on $|w|\le2\eta$
\[
 L_\eta=O(\eta),\qquad
 M_\eta=M_0(s,0)+O(\eta),\qquad N_\eta=O(1).
\]
For example the $ss$ derivative of the added normal displacement
is $O(\eta^2)$, and the change of normal in the bounded original
$ss$ coefficient contributes $O(\eta)$; the mixed and $ww$
estimates follow in the same way. Hence
\[
 L_\eta N_\eta-M_\eta^2
     =-M_0(s,0)^2+O(\eta)<0
\]
uniformly for small $\eta$, since $\min_s|M_0(s,0)|>0$.
The first fundamental determinant stays bounded below by a
positive constant by $C^1$ convergence. Formula
\eqref{eq:gauss-curvature} proves $K_{X_\eta}<0$ on the
modified collar. Outside it the surface is unchanged, so the
curvature is negative on all of $U$.

Fix $\eta$ small enough for all the preceding conclusions and for
$\|X_\eta-X_0\|_{C^1}<\varepsilon$.
On $|w|<\eta$, $X_\eta=X_0+hn=Z\circ q^{-1}$ exactly.
By compactness of the core, there is $a>0$ such that the image
under $q$ of $|u|<a$ lies in $|w|<\eta$. Restricting $q$ to
this strip gives exactly the inclusion asserted in the statement.
If $X_0$ is tight, Lemma~\ref{lem:tight-preservation} applies.
\end{proof}

\begin{lemma}[A supported band gives a global bending]
\label{lem:globalize}
Let $X:M\to\R^3$ be an embedding of a closed surface, let
$\ell>0$, let $I\subset\R$ be a nonempty open interval, and put
$A=(\R/\ell\mathbb Z)\times I$. Let $\iota:A\to M$ be a smooth
annular embedding, and set $B=X\circ\iota$.
If $B$ has an infinitesimal bending $Y_A$ compactly supported in $A$
and with $\int_A\det(Y_A,B_s,B_c)\,ds\,dc\ne0$, then $X$ has an
infinitesimal bending $Y$ supported in $\iota(\supp Y_A)$ and
with nonzero outward normal flux.
\end{lemma}
\begin{proof}
An embedding of equal-dimensional manifolds has open image, by
the inverse function theorem. Define $Y\circ\iota=Y_A$ on that
image and $Y=0$ outside it.
The image of $\supp Y_A$ is a compact subset of the open set
$\iota(A)$. Every point outside $\iota(A)$ therefore has a
neighborhood disjoint from this image, on which the extended
field is zero. Thus the extension is smooth.
The chain rule and the identity $X\circ\iota=B$ transfer
\eqref{eq:strain} to $\iota(A)$; outside the support it is zero.
On this connected annulus the outward normal of $X$ is either
the parameter normal $(B_s\times B_c)/|B_s\times B_c|$
or its negative, with a single constant sign. The
change-of-variables formula therefore gives
\[
 \int_M\langle Y,\nu_X\rangle\,dA_X
      =\pm\int_A\det(Y_A,B_s,B_c)\,ds\,dc\ne0.
\]
\end{proof}

\begin{corollary}[One curve implies a supported bending]\label{thm:one-curve}
Suppose $X_0:T^2\to\R^3$ is a smooth tight embedding and
$c:S^1\to T^2$ is a smooth embedding such that $X_0\circ c$ is
asymptotic and $K_{X_0}<0$ along $c(S^1)$.
For every open set $U$ with
\[
 c(S^1)\subset U\subset\{K_{X_0}<0\},
\]
and every $\varepsilon>0$, there exist a smooth tight embedding
$X:T^2\to\R^3$ and an infinitesimal bending $Y$ of $X$ such that
\begin{equation}\label{eq:one-curve-output}
 \begin{gathered}
 \supp(X-X_0)\Subset U,\qquad K_X<0\ \hbox{on }U,\qquad
 \|X-X_0\|_{C^1}<\varepsilon,\\
 \supp Y\Subset U,\qquad
 \int_{T^2}\langle Y,\nu_X\rangle\,dA_X\ne0.
 \end{gathered}
\end{equation}
\end{corollary}

\begin{proof}[Proof of Corollary~\ref{thm:one-curve}]
Reparametrize the original core and its source embedding $c$ by
arclength. Choose a collar with compact closure in $U$.
Proposition~\ref{prop:inflection} supplies an inflection on the
original curve. Apply Proposition~\ref{prop:prepare}, with
$C^2$ error as small as needed, supported in that collar.
Call the resulting tight embedding $X_1$.
Reparametrize its new core and the source embedding $c$ by
arclength, and equip the curve with its Darboux frame. Write $s$ and
$\ell$ for this arclength and its period. The curve satisfies the
hypotheses of Proposition~\ref{prop:ruled}.

Fix the ruling supplied by that proposition.
Apply Theorem~\ref{prop:insert} to $X_1$ with $Z=R$,
again with support compactly in $U$, to obtain $X$.
Choose the errors in the two modifications so that their sum
in $C^1$ is less than the prescribed $\varepsilon$.
The surface remains negatively curved on all of $U$ and is tight.
The inclusion of a neighborhood of the ruled core is exact.

Let $q$ and $a>0$ be supplied by Theorem~\ref{prop:insert}.
Only now choose a bounded nonempty open interval $I$ of sufficiently
large leaf parameters $c$ so that $c+\omega(s)>0$ and
$0<u(s,c)<a$ for every $(s,c)\in(\R/\ell\mathbb Z)\times I$.
The annulus
$B(s,c)=\gamma(s)+\rho(s)V(s)/(c+\omega(s))$
lies inside that exact neighborhood.
Set $\iota(s,c)=q(s,u(s,c))$.
Then $\iota$ is an annular embedding and
$X\circ\iota=B$. Its constant-$c$ circles are closed and
asymptotic by Proposition~\ref{prop:ruled}; its constant-$s$
curves are straight, and its curvature is negative by
\eqref{eq:ruled-curvature}.
Thus Theorem~\ref{thm:band-bending} supplies a field
on this annulus with compact support and nonzero signed
flux. Lemma~\ref{lem:globalize} produces the desired global
field $Y$. Its compact support lies in $U$.

This proves \eqref{eq:one-curve-output}. The ruling is fixed before
the insertion width, the leaf interval, and the bending profile
are chosen.
\end{proof}

\section*{Acknowledgements}
I thank my mentor Xiang Ma for helpful discussions and encouragement.
AI-assisted tools were used in the preparation of this manuscript.

\begin{thebibliography}{99}
\bibitem{Aleksandrov}
V.~A.~Aleksandrov,
\emph{Sabitov's conjecture that volume is stationary under infinitesimal
bending of a surface}, Siberian Math. J. \textbf{30} (1989), 678--684,
\href{https://doi.org/10.1007/BF00971257}{doi:10.1007/BF00971257}.
\bibitem{BK}
T.~F.~Banchoff and W.~K\"uhnel,
\emph{Tight submanifolds, smooth and polyhedral}, in
\emph{Tight and Taut Submanifolds}, MSRI Publications \textbf{32},
Cambridge University Press, 1997, 51--118,
\href{https://library.slmath.org/books/Book32/files/banchoff.pdf}{available from SLMath}.
\bibitem{BGQR}
H.~Banks, M.~Ghomi, H.~Qiu, and M.~Raffaelli,
\emph{A counterexample to Nirenberg's conjecture on closed asymptotic
curves}, \href{https://arxiv.org/abs/2610.13848v2}{arXiv:2610.13848v2},
21 October 2026.
\bibitem{doCarmo}
M.~P.~do Carmo, \emph{Differential Geometry of Curves and Surfaces},
Prentice-Hall, 1976.
\bibitem{Fenchel}
W.~Fenchel, \emph{\"Uber Kr\"ummung und Windung geschlossener Raumkurven},
Math. Ann. \textbf{101} (1929), 238--252,
\href{https://doi.org/10.1007/BF01454836}{doi:10.1007/BF01454836}.
\bibitem{GR}
M.~Ghomi and M.~Raffaelli,
\emph{Topology of closed asymptotic curves on negatively curved surfaces},
J. Geom. Anal. \textbf{35} (2025), article 381,
\href{https://doi.org/10.1007/s12220-025-02200-3}{doi:10.1007/s12220-025-02200-3};
\href{https://arxiv.org/abs/2412.19266v2}{arXiv:2412.19266v2}.
\bibitem{Hirsch}
M.~W.~Hirsch, \emph{Differential Topology},
Graduate Texts in Mathematics \textbf{33}, Springer, 1976,
\href{https://doi.org/10.1007/978-1-4684-9449-5}{doi:10.1007/978-1-4684-9449-5}.
\bibitem{Kovaleva}
G.~A.~Kovaleva,
\emph{Absence of closed asymptotic lines on tubes of negative Gaussian
curvature with a one-to-one spherical mapping},
Fundam. Prikl. Mat. \textbf{1} (1995), no.~4, 953--977 (Russian).
\bibitem{Lee}
J.~M.~Lee, \emph{Introduction to Smooth Manifolds}, second edition,
Graduate Texts in Mathematics \textbf{218}, Springer, 2013,
\href{https://doi.org/10.1007/978-1-4419-9982-5}{doi:10.1007/978-1-4419-9982-5}.
\bibitem{Spivak}
M.~Spivak, \emph{A Comprehensive Introduction to Differential Geometry},
Vol.~V, third edition, Publish or Perish, 1999, Chapter~12.
\bibitem{Teschl}
G.~Teschl, \emph{Ordinary Differential Equations and Dynamical Systems},
Graduate Studies in Mathematics \textbf{140}, American Mathematical
Society, 2012,
\href{https://www.mat.univie.ac.at/~gerald/ftp/book-ode/ode.pdf}{author-hosted text}.
\end{thebibliography}
\end{document}


